% !PRINTED 696
\weberpage{696}
A function $z$ which becomes 0 and $\infty$ at only a single point each,
and so has order 1, therefore exists only in the case $p = 0$.
Every function in $\Omega$ \gesperrt{can be rationally represented by such a $z$}%
\footnote{This function field leads to what the geometers call
``\gesperrt{unicursal curves}''.}.

If $r_i > 2$, then by § 198 $\mu_i$ and $z\mu_i$ are differential
quotients of the first kind. By (1) there are therefore polygons $\mathfrak{W}$, $\mathfrak{W}'$
of the first kind which satisfy the condition:
\[
\mu_i = \frac{\mathfrak{A}^2\mathfrak{W}}{\mathfrak{Z}},\quad \mu_i z = \frac{\mathfrak{A}^2\mathfrak{W}'}{\mathfrak{Z}},\quad z = \frac{\mathfrak{W}'}{\mathfrak{W}} = \frac{\mathfrak{A}'}{\mathfrak{A}},
\]
and since $\mathfrak{A}$, $\mathfrak{A}'$ are coprime, $\mathfrak{A}$ must divide $\mathfrak{W}$, and $\mathfrak{A}'$ must divide $\mathfrak{W}'$.
Hence $A$ is a class of the first kind.

If $A$ is a class of the second kind, then all the exponents
$r_1, r_2, \ldots, r_n$ are therefore equal to 2 or less than 2, and in
particular $r_{s+1}=2, r_{s+2}=2, \ldots, r_n=2$. Consequently
\[
p = (r_2-1)+(r_3-1)+\cdots+(r_n-1) = n-s,
\]
and by (2) we have the theorem:

3. \gesperrt{The dimension of a proper class $A$ of
the second kind is}

\weq{3}{(O,A) = n - p + 1.}

If, further, $A$ is a proper class of the first kind and

\weq{4}{AB = W}

is the complete class of the first kind, then take a polygon $\mathfrak{B}$
in $B$ and form the differential quotients of the first kind:
\[
w = \frac{\mathfrak{A}^3\mathfrak{B}}{\mathfrak{Z}},\quad zw = \frac{\mathfrak{A}^2\mathfrak{A}'\mathfrak{B}}{\mathfrak{Z}},
\]
and this form is preserved when $\mathfrak{B}$ is replaced by an equivalent
polygon $\mathfrak{B}'$. The dimension $(O, B)$ is therefore as great
as the dimension of the family of differential quotients of the first
kind $w$ which have the property that $zw$ too is still of
the first kind. For this to be the case, in the
expression:
\[
w = y_1\mu_1 + y_2\mu_2 + \cdots + y_n\mu_n
\]
the degrees of the functions $y_1, y_2, \ldots, y_n$ may rise only to the height
$r_1-2, r_2-2, \ldots, r_n-2$, and from this it follows:
% [sic] printed r_1-2, r_2-2, ..., r_n-2 as the height to which the degrees of y_1, ..., y_n may rise, but that is the bound for w alone to be of the first kind (cf. r_i > 2 above); for zw to be of the first kind as well the degrees can reach only r_i-3, which is what the count (O,B) = (r_{s+1}-2)+...+(r_n-2) on p. 697 presupposes
%CONT>>
